\egasection{§8\quad 例外超平面集合的维数}\source{\texttt{EGAVaiello.pdf} 第21--36页；Dot各批连续译稿并入}
\egap{8.1}
在前面各节，尤其是第2节和第3节，我们已给出一些断言：使得 $Y$ 具有某个性质 $P$ 的 $\xi\in P^\vee$ 所成的集合是可构造的，且包含泛点 $\eta$；或者说，对性质 $P$ 而言“例外”的 $\xi\in P^\vee$ 所成的集合 $Z_P$ 是可构造且稀疏的，即其闭包的余维至少为一。〔注意：这里假设 $S=\Spec(k)$。〕
\transnote{底本首句重复了“使得……的 $\xi\in P^\vee$ 所成的集合”这段引导语。本译文只排一次，在此记录重复。}
在某些情形，可以给出这个余维的一个更好的下界，从而使该断言更为精确；这对于某些问题很重要。例如，如果看出这个余维大于或等于二，那么 $P^\vee$ 中一条“充分一般”的直线 $D$ 就不与 $Z_P$ 相交，因而存在（若 $k$ 无限）超平面截面 $Y_\xi$ 的“线性束”（$\xi$ 为 $D$ 的几何点），其中所有截面都具有性质 $P$；例子见第〔节号缺失〕节。\ednote{节号遗漏，请询问 A.G.。〔底本脚注29。〕}
\transnote{底本为 upper bound（上界），但本段要把“余维至少一”加强为“余维至少二”；正确的方向是改善下界，正文据此订正。}
从编写的角度看，本节的结果使前几节的结果更精确，因而产生一个问题：是否应当把这些补充另列一节（或一编号），还是在前文逐步推进时就给出更精确的版本？\textit{Redactor decidetur}〔拉丁文〕。\ednote{由编辑决定。〔底本脚注30。〕}

\egap{8.2}
令 $Z$ 为满足 $\dim Y_\xi>\dim X-1$ 的 $\xi\in P^\vee$ 所成的集合。假设对 $X$ 的每个不可约分支 $X_i$，都有
\[
\dim f(X_i)>\text{〔字迹不清：是二吗？请询问 A.G.〕}\text{。}
\]
那么 $Z$ 在 $P^\vee$ 中的余维至少为二。这由（2.1）、（2.2）（后者蕴含 $Y$ 的每个不可约分支都支配 $P^\vee$）以及态射 $Y\to P^\vee$ 的维数理论得出。从这个结果出发，可以把下述情形作为推论：从 $X$ 的一个闭子集 $Z$ 出发，考察其在 $Y_\xi$ 中的逆像 $Z_\xi$ 的维数（$\xi\in P^\vee$）；甚至可以考察如下 $\xi\in P^\vee$ 的集合：给定闭子集 $T$ 的基变换 $T_{k(\xi)}$ 存在一个不可约分支，其在 $Y_\xi$ 上的迹具有最大的维数。〔注意：始终假设对 $Z$ 的每个不可约分支 $Z_i$ 有 $\dim f(Z_i)>0$。〕

最后，这个方向上最精确的断言，是如下修改后的断言；它容易由第一个断言（$X$ 不可约的情形）和（2.7）推出：设 $F$ 是 $X$ 上的凝聚模，假设对 $F$ 的每个相伴素循环 $T$ 都有 $\dim f(T)>0$，则使 $\phi_\xi$ 不是 $F_{k(\xi)}$-正则的 $\xi\in P^\vee$ 所成的集合是（可构造的，并且）余维至少为二。（$\phi_\xi$ 的记号来自第5节。）可以将此作为主要断言，把前面的断言列作推论，而证明则由其中一个推论给出或经由它进行。

请注意，沿用上述记号，若 $\xi\in P^\vee-Z$，则对每个 $y\in Y_\xi$ 都有
\[
\coprof_y F_{k(\xi)}=\coprof_y G_\xi\text{。}
\]
因此，若 $\coprof F\leq n$，则对 $\xi\in P^\vee-Z$，有 $G_\xi$ 的余深度 $\leq n$；特别地，若 $F$ 是 Cohen--Macaulay 的，则对 $\xi\in P^\vee-Z$，$G_\xi$ 也是 Cohen--Macaulay 的。最后，若 $F$ 满足 $(S_k)$，则对 $\xi\in P^\vee-Z$，$G_\xi$ 满足 $(S_{k-1})$（参见 $0_{\mathrm{IV}}$）。
\transnote{底本把给定闭子集与参数例外集均记为 $Z$，中间又出现 $T_{k(\xi)}$；正文将后者明确为给定闭子集 $T$ 的基变换，并不改变所取例外集。余深度公式左端原作 $G_{k(\xi)}$，下一句又作 $T_\xi$；按 $G_\xi=F_{k(\xi)}/\phi_\xi F_{k(\xi)}$ 的正则商关系，分别订为 $F_{k(\xi)}$、$G_\xi$：正则元使局部维数与深度同时降一，故余深度不变。原文“余维为二”过强，例如 $X=P=\mathbf P^3$ 的恒等嵌入具有空异常集；正文仅写所需下界。开头那个难辨的维数阈值仍未补字。}

\egap{8.3}
我们注意到：如果 $F$ 满足 $(S_k)$，对于某个 $\xi\in P^\vee$，使得 $\phi_\xi$ 是 $F_{k(\xi)}$-正则的，并且 $G_\xi$ 有一个余维至少为二的分支\ednote{字迹不清，请询问 A.G.。〔底本脚注31。〕}，即使 $F=\mathcal O_X$、$k=1$，且 $X$ 几何整、维数为二；或者（$k=2$，$X$ 几何整且几何正规、维数为三）。〔此句按底本保留残缺结构。〕
只须从代数闭域 $k$ 上的一张射影整曲面
\[
X\subset P^r
\]
出发，它有一点 $x$，使 $X$ 在该点不是 Cohen--Macaulay 的；那么，对每个经过 $x$ 的超平面，相应的超平面截面 $Y_\xi$ 都以 $x$ 为一个嵌入相伴循环。相应地，从三维正规（因此满足 $(S_2)$）整簇 $X\subset P$ 出发，它有一点 $x\in X$，使 $X$ 在此处不是 Cohen--Macaulay 的；则经过 $x$ 的那些 $Y$ 不是 CM 的，也就是说，它们在 $x$ 处不满足 $(S_2)$。

在这些例子中，对性质 $(S_k)$ 而言“例外”的 $\xi$ 所成的集合包含由 $x\in P$ 定义的 $P^\vee$ 中的超平面，因而其余维为一（而非至少二）。关于这一方向的一般精确结果，比较下面的（8.5）？
\transnote{首句在底本中缺少完整谓语，仍按残句保留。其中 $f_{k(\xi)}$ 订为模 $F_{k(\xi)}$；后面的环境 $P^\vee$、点 $X\in X$ 分别订为 $P$、$x\in X$。“met $(S_2)$”不成英文，且这里正规三维局部环非 Cohen--Macaulay 时深度为2，正则超平面商的维数为2而深度为1，故应为“不满足 $(S_2)$”。这些可确定订正不恢复首句所缺的完整断言。}

\egas{8.4}{命题}
设 $T$ 是 $X$ 的闭子集，并假设 $\codim(T,X)\geq k$。那么，对每个 $\xi\in P^\vee$，都有 $\codim(T_\xi,Y_\xi)\geq k-1$。令 $Z$ 为满足
\[
\codim(T_\xi,Y_\xi)=k-1
\quad\text{（即 }\codim(T_\xi,Y_\xi)<k\text{）}
\]
的 $\xi\in P^\vee$ 所成的集合，则 $Z$ 是 $P^\vee$ 的可构造、无处稠密〔法文：稀疏〕子集，即 $\overline Z$ 在 $P^\vee$ 中的余维至少为一。

它的余维至少为二的充要条件是：对 $T$ 的每个在 $X$ 中余维等于 $k$ 且满足 $\dim f(T_i)=0$ 的不可约分支 $T_i$，都存在 $X$ 的一个不可约分支 $X_j$，使
\[
\codim(T_i,X_j)=k,\qquad \dim f(X_j)=0.
\]
〔底本在此附有“或者‘闭包’被划去了？”的字迹疑问。〕即，若 $f$ 拟有限且 $k>0$，则 $T$ 没有满足 $\dim_xX=k$ 的孤立点 $x$。\cmt{D17-E02}{底本把 $T_i$ 写成 $X$ 的分支，且写 $\dim f(X_j)-0$、$\dim_xX-k$；正文依同条的证明和随后的引理，分别订为 $T$ 的分支及两个等号。对 $k>0$，$X$ 的不可约分支本身在 $X$ 中余维为0，不可能符合原句要求；等号则是后文反复使用的条件。完整旧译及原疑字记录保存在报告。}
第一个断言立即由下面的引理（8.4.1）(a)得出；这个引理是对第5节的一项事后补充。

\egas{8.4.1}{引理}
设 $X$ 为局部诺特预概形，$L$ 为 $X$ 上的可逆模，$\phi$ 为 $L$ 的一个截面，$Y=V(\phi)$，$T$ 为 $X$ 的闭子集。假设
\[
\codim(T,X)\geq k.
\]
则：
\begin{enumerate}
\item[(a)] $\codim(T\cap Y,Y)\geq k-1$。
\item[(b)] 要有
\[
\codim(T\cap Y,Y)=k-1,
\quad\text{即 }\codim(T\cap Y,Y)<k,
\]
其充要条件是：存在 $T$ 的一个包含于 $Y$ 的不可约分支 $T_i$，使
\[\codim(T_i,X)=k,\]
并且对 $X$ 的每个包含 $T_i$ 且满足
\[
\dim\mathcal O_{X_j,T_i}=\dim\mathcal O_{X,T_i}\quad(=k)
\]
的不可约分支 $X_j$，都有 $X_j\not\subset Y$。
\end{enumerate}
根据 $0_{\mathrm{IV}}$、第IV章中关于维数的一般事实，这个引理的验证是直接的。\cmt{D17-E02续}{底本假设印作 $\codim(Y,X)\geq k$，按字面不成立：令 $X=\mathbb A^1$、$Y=\{0\}$、$T=X$、$k=1$，则左侧等式成立，却不存在所要求的 $T_i\subset Y$。正文将该假设订为本命题所需的 $\codim(T,X)\geq k$；报告给出局部维数论证明，分别验证(a)与(b)，故此处不只是凭上下文猜换字母。原作者“验证直接”的详略保持，补证另行标明为译校所加。}

在（8.4）的假设下，由（8.4.1）(b)可以看出哪些超平面 $H_\xi$ 是例外的。排除如下 $\xi\in P^\vee$ 所成的集合 $Z_0$：$T$ 或 $X$ 有一个不可约分支 $R$，满足 $\dim f(R)>0$，且 $R_\xi$ 的“维数过大”。这是一个余维至少为二的集合，在以下论述中不予计入。例外的 $H_\xi$ 是那些使得存在一个 $T_i$ 满足
\[
\codim(T_i,X)=k,\qquad\dim f(T_i)=0,\qquad f(T_i)\subset H_\xi,
\]
并且，对 $X$ 的每个满足 $X_j\supset T_i$ 及 $\codim(T_i,X_j)=k$ 的不可约分支 $X_j$，都有 $f(X_j)\not\subset H_\xi$ 的超平面。\transnote{底本写 $\dim f(T)=0$、$f(T)\ \mathrm{CH}$，其中英译编者脚注32说“大概是 $H_\xi$”。按引理(b)正在考察的分支及其包含条件，正文明确写 $T_i$ 和 $f(T_i)\subset H_\xi$；不是从字形推断未辨的“CH”。$Z_0$ 的余维只需至少为二，不要求精确等于二。}

给定满足 $\codim(T_i,X)=k$ 的 $T_i$，若存在包含它的 $X_j$，使 $\codim(T_i,X_j)=k$ 且 $\dim f(X_j)=0$，则 $f(X_j)=f(T_i)$；所以，只要 $H_\xi$ 包含 $f(T_i)$，它就同时包含 $f(X_j)$，因而 $\xi$ 相对于 $T_i$ 不是例外的。反之，如果对每个包含 $T_i$ 且满足 $\codim(T_i,X_j)=k$ 的 $X_j$ 都有 $\dim f(X_j)>0$，则对 $\xi\in P^\vee-Z_0$，$\xi$ 相对于 $T_i$ 为例外的充要条件是 $f(T_i)\subset H_\xi$；这样的 $\xi$ 所成的集合，是 $P^\vee$ 中一个超平面在 $P^\vee-Z_0$ 上的迹。\transnote{底本在第一处余维等式的右端为问号，并附“字迹不清，请询问Grothendieck”，又把两个包含关系写成“不包含”，把 $\dim f(X_j)>0$ 写成 $f(X_j)>0$，把参数开集写成 $P-Z_0$。正文按引理(b)明确所考察的是余维恰为 $k$ 的分支，并订正各关系。固定点不在 $H_\xi$ 中的参数构成开集；包含该点的超平面才组成对偶超平面，故末处必须是包含。}
这就证明了（8.4），并且证明了更精确的结果：在除去一个余维至少为二的集合后，例外集正是若干超平面在该开集上的迹；这些超平面可由以上证明明显地确定。\cmt{D17-E02再续}{底本把结论写成“例外集是一个余维至少为二的集合与若干超平面的并”。一般只能断言在 $P^\vee-Z_0$ 上的上述等式。例如令 $X$ 是 $\mathbf P^2$ 中的一条直线，$T=\{x\}$，$k=1$；异常参数是经过 $x$ 的所有直线中除去 $X$ 本身，因而为一条对偶直线删去一点，并不是若干整条超平面与低维集合的并。正文保留证明真正给出的开集上的等式。}

（恐怕这段写法很松散，或者说很潦草〔英译者措辞说明〕，因为我一直在几何意义下推理，却没有说出来，即取值于代数闭域的点。当然，（8.4）中所陈述的条件确实是几何的，所以可以假设 $k$ 代数闭，并对 $k$-有理点进行论证。）利用（8.4）、（5.7.4）以及（8.2）的末尾，得到：

\egas{8.5}{推论残文（首项假设及 $Z_n$ 定义缺失）}
假设对所有相伴素循环 $R$，我们有“至多仅仅”〔字迹不清〕\ednote{字迹不清，请询问 A.G.。〔底本脚注33。〕}，并假设 $F$ 满足 $(S_k)$。

要使 $P^\vee$ 中满足 $\phi_\xi$ 为 $F_{k(\xi)}$-正则且 $G_\xi$ 满足 $(S_k)$ 的点所成的（可构造）集合的补集余维至少为二，其充要条件如下：对每个整数 $n\geq0$，以 $Z_n$ 表示满足 $\coprof_x$〔字迹不清〕\ednote{请询问 A.G.。〔底本脚注34。〕}的 $x\in T=\operatorname{Supp} F$ 所成的集合。对 $Z_n$ 的每个不可约分支 $Z_{ni}$，若
\[
\codim(Z_{ni},T)=n+k+1,\qquad\dim f(Z_{ni})=0,
\]
则存在 $T$ 的一个包含 $Z_{ni}$ 的不可约分支 $T_j$，使
\[
\codim(Z_{ni},T_j)=n+k+1,\qquad\dim f(T_j)=0\text{。}
\]

当 $f$ 拟有限时，对〔字迹不清，请询问 A.G.〕的每个闭子集 $R$ 都有 $\dim f(R)=\dim R$，所以判据成为：任一 $Z_n$ 中都不存在使 $\dim_zT$（$=\dim F_z$）等于 $n+k+1$ 的孤立点 $z$。当 $F$ 等维且维数为 $d$ 时，若 $d\leq k$，这个条件是空的（实际上我们早已知道，因为这时 $F$ 的 $(S_k)$〔假设〕恰是 Cohen--Macaulay 假设）；若 $d\geq k+1$，其意思是，$T$ 中 $F$ 的余深度 $>d-(k+1)$ 的点所成的集合 $Z_{d-(k+1)}$，也就是说这里对应的是 $F$ 的真深度 $<k+1$（尽管事先由 $(S_k)$ 和 $k\leq d$ 只能得到 $F$ 的真深度 $\geq k$）。〔原句在此结束，缺少完整谓语。〕如果不再假设 $F$ 等维，仍可用下面的简单方式表达所需条件：
\transnote{$Z_n$ 的定义和首项假设均有实质缺文，不可据此直接使用判据。底本“余深度 $>d-(k+1)$”后写成“真深度 $\geq k+1$”，方向相反；由余深度等于局部维数减深度，在维数 $d$ 处应为深度 $<k+1$，正文已订正。缺少的谓语和完整判据不据此猜补。}

\egap{8.6}\textbf{〔承接前条的判据残文；下列明确错符已校正，完整前提仍待核。〕}
对每个满足 $\dim F_x\geq k+1$ 的闭点 $x\in\operatorname{Supp} F$，都有 $\prof F_x\geq k+1$。

取 $Z=x$，立即可见其充分性。必要性由下述观察得出：对于使 $\phi_\xi$ 为 $F_{k(\xi)}$-正则且 $x\in Y_\xi$ 的每个 $\xi$，都有
\[
\dim G_{\xi,x}=\dim F_x-1,\qquad
\prof G_{\xi,x}=\prof F_x-1\text{。}
\]
所以，若 $x$ 使上述条件失效，就有 $\prof G_{\xi,x}<k$，但 $\dim G_{\xi,x}\geq k$，这表明 $G_\xi$ 在 $x$ 不满足 $(S_k)$；然而，使 $x\in Y_\xi$ 的 $\xi$ 所成的集合余维为一。〔注意：我隐含地假设了 $k$ 代数闭，而我们立即可以归结到这种情形。〕前述一般判据在（8.6）的情形应当是显然的。
\transnote{底本首行作 $k_1$，但此段考察的正是维数至少 $k+1$ 而深度未达到 $k+1$ 的点；按随后维数降一、深度降一的论证订为 $k+1$。$X_{k(\xi)}$、$\prof F)_x$、$G_xix$ 分别订为 $F_{k(\xi)}$、$\prof F_x$、$G_{\xi,x}$。底本末段作 $\prof G_xix\geq k$；这与其声称不满足 $(S_k)$ 相反，且若 $F$ 满足 $(S_k)$、$\dim F_x\geq k+1$ 而此处深度条件失效，就有 $\prof F_x=k$，故 $\prof G_{\xi,x}=k-1<k$。以上为可确定的推演修正，不补全（8.5）已缺失的假设及 $Z_n$ 定义。}

% DOT_IMPORT 03 8.6后预备论证_校订续译工程

\subsection*{余法丛与奇异零点}
\noindent\textbf{译校记号约定。}沿用第1节 $X/S$ 有限型的背景。令 $P=\mathbf P(E)$，$P^\vee=\mathbf P(E^\vee)$，采用一维商空间的射影化约定。令 $W=X\times_S P^\vee$。在相对维数恒定的开集上，记 $d=\dim(X/S)$、$r=\dim(P/S)=\dim(P^\vee/S)$；一般情形按这样的开集局部解释。记 $N^\vee$ 为余法模，$N=(N^\vee)^\vee$ 为法模。$\mathcal O_W(a,b)$ 中两项分别来自 $\mathcal O_X(1)=f^*\mathcal O_P(1)$ 和 $\mathcal O_{P^\vee}(1)$。同样的记号用于其在 $Y$ 或测试预概形 $Z$ 上的拉回。底本的 $\nu_{X/P}$ 对应这里的 $N$。


我们现在研究 $Y$ 中那些点 $y$：在这些点处，相应的 $Y_\xi$ 相对于 $k(\xi)$ 不光滑。限于 $f:X\to P$ 非分歧（实际应用中它将是一个浸入）且 $X\to S$ 光滑的情形。不必假设 $S$ 是域的谱。由于 $f$ 非分歧，典范同态
\[
f^*\Omega^1_{P/S}\longrightarrow\Omega^1_{X/S}
\]
是满射，其核是 $X$ 上的局部自由模，记为 $N^\vee$。当 $f$ 是浸入时，在把 $X$ 视为某个开集中的闭子概形后，它就是由 $X$ 的理想 $J$ 定义的余法模 $J/J^2$；在一般情形也把它称为余法模。\transnote{底本(a)中间项印成 $f^*\Omega^1_{P/X}$；正文改为 $f^*\Omega^1_{P/S}$。已给定的是 $X\to P\to S$，且同段先前已写出这条相对微分满射。$X/S$ 光滑使右端局部自由，故该满射局部分裂，核局部自由。}
\begin{equation}
0\longrightarrow N^\vee\longrightarrow f^*\Omega^1_{P/S}
\longrightarrow\Omega^1_{X/S}\longrightarrow0.\tag{a}
\end{equation}
注意，在 $P$ 上还有一条典范正合列（例如应在第16节中作为例子列出）：
\begin{equation}
0\longrightarrow\Omega^1_{P/S}(1)\longrightarrow E_P
\longrightarrow\mathcal O_P(1)\longrightarrow0.\tag{b}
\end{equation}
也就是说，$\Omega^1_{P/S}$ 典范同构于由 $E_P\to\mathcal O_P(1)$ 导出的典范同态 $E_P(-1)\to\mathcal O_P$ 的核。对它施加 $f^*$，得到
\begin{equation}
0\longrightarrow f^*(\Omega^1_{P/S}(1))\longrightarrow E_X
\longrightarrow\mathcal O_X(1)\longrightarrow0.\tag{$b_1$}
\end{equation}
这在 $X$ 上给出了 $f^*\Omega^1_{P/S}\otimes\mathcal O_X(1)$ 的明确描述，因此可将 $N^\vee(1)$ 视为 $E_X$ 的一个局部直因子子模；等价地，其对偶 $N(-1)$ 典范同构于 $E_X^\vee$ 的一个商模。因而
\[
\mathbf P(N(-1))\simeq\mathbf P(N)
\]
可以典范地嵌入
\[
\mathbf P(E_X^\vee)=X\times_S P^\vee=W,
\]
作为 $X$ 上的射影子丛，从而作为闭子概形。它必定包含于 $Y$，因为 $N^\vee(1)$ 包含在 $E_X\to\mathcal O_X(1)$ 的核中。\transnote{底本此处将这个子模写成 $\Omega_{X/P}(1)$；但 $f$ 非分歧时 $\Omega^1_{X/P}=0$，实际用于构造的子模是上文已嵌入 $E_X$ 的 $N^\vee(1)$。底本紧接着附有“最后两个符号难辨，请询问 A.G.”的说明。此处译校以两条正合列确定该包含关系。}

这个预概形的底空间恰是 $Y=V(\phi)$ 中那些点所成的集合：它们是 $\mathcal O_W(1,1)$ 的截面 $\phi$ 相对于基底 $P^\vee$ 的奇异零点（第16节）。\ednote{参见这些笔记的第二部分。〔底本脚注35。〕}也就是说，对一个在 $P^\vee$ 上取值的域 $K$，它的 $K$-值点就是 $Y_K\subset X_K$ 中使 $\phi_K$ 在该点至少二阶消失的点；等价地，$Y_K$ 在那里不是相对维数 $d-1$ 的光滑 $K$-预概形。\transnote{底本这里写 $r-1$，下一页相应的超平面截面条件却写 $d-1$；同一论证末尾明定 $r$ 为环境射影空间的相对维数。此处译校以 $d$ 表示 $X/S$ 的相对维数，截面预期维数应为 $d-1$。另用 $K$ 表示此处的取值域，以免与稍后 $S=\Spec(k)$ 混淆。要求的是“光滑且具有预期维数”，不能删去维数条件。}
上面所述对奇异零点的刻画，即把它们看作 $W$ 中光滑子概形 $\mathbf P(N)$ 的点，立即给出下述断言，值得作为预备命题列出。\transnote{底本在“刻画”附近留有一处“字迹不清，请询问 A.G.”标记。这里仅按前后两个完整句之间已经明确的数学关系衔接，不声称复原了难辨原字。这里的光滑是相对于 $S$，不是声称这个闭浸入是光滑态射。}

若 $S=\Spec(k)$，$H$ 为 $P$ 的超平面，则 $Y=X\times_P H$ 在点 $x\in Y(k)$ 处相对于 $k$ 光滑且相对维数为 $d-1$（即 $x$ 是 $H$ 所定义的 $\mathcal O_X(1)$ 的截面 $\phi$ 的非奇异零点，亦即几何非奇异零点），当且仅当 $H$ 不包含 $X$ 在 $x$ 处相对于 $k$ 的切空间经 $f$ 诱导的像所对应的射影线性子空间。\transnote{底本印作“经 $\phi$ 的像”。$\phi$ 是线丛截面，不能在此扮演 $X\to P$ 的角色；应由 $f$ 的切映射以及下一段的线性子空间对应来解释。}换句话说，当 $f:X\to P$ 是浸入，从而可将 $X$ 与 $P$ 的一个子概形相认时，条件就是：$H$ 在 $x$ 处不与 $X$ 相切。

一旦把这一陈述的含义说清楚，即说明 $P$ 在点 $a=f(x)$ 处的切空间的一个向量子空间怎样确定 $P$ 的一个线性子空间，使得“$H$ 不包含该子空间”这句话有意义，该断言便直接由光滑性的 Jacobian 判据或奇异零点的定义得出。当然，这来自上面的正合列(b)：它给出 $a$ 处切空间的向量子空间与 $P$ 中包含 $a$ 的线性子空间之间的一一对应。\transnote{底本写“factor subspaces”，在域上每个向量子空间都有补空间。本译文按其紧接的解释译为向量子空间，不改译为商空间。}无论如何，这个对应就是把经过 $a$ 的线性子簇对应到它在 $a$ 处的切空间，并将后者看作 $P$ 在 $a$ 处切空间的子空间。

这些“连锁事实”，连同有关线性子簇及 Grassmann 簇的各种类似事实，应该集中在一两个预备编号或预备段落中，当然要对任意基底陈述。事实上，我们能做得更好：第16节定义的 $\phi$ 相对于 $P^\vee$ 的奇异零点预概形 $Y^{\mathrm{sing}}$，恰是 $\mathbf P(N)$。\transnote{底本第26页此处排作 $\mathbf P(\nu_{X/P}^{\vee}$，还缺少右括号。应为 $\mathbf P(N(-1))\simeq\mathbf P(N)$：一维商的射影化由 $E_X^\vee\twoheadrightarrow N(-1)$ 给出，不能在此把法模改成余法模。}它相对于 $S$ 光滑；在非空的相关等维部分，其相对维数为\transnote{当 $N$ 的秩为零时，其射影化为空；此处译校明确非空部分，以免把空概形错误地断言为实际维数等于 $r-1$。$N$ 的秩为 $r-d$，所以非空纤维为 $\mathbf P^{r-d-1}$。}
\[
d+(r-d-1)=r-1.
\]
因此我们处于第16号或第16节所研究的有利情形。\ednote{请询问 A.G. 这个引文的位置：应在这些笔记稍后的部分。〔底本脚注36。〕}为了验证上述断言，注意按定义，$Y^{\mathrm{sing}}$ 正是下述截面的零点在 $Y$ 中定义的子预概形：
\[
\Psi=(d_{W/P^\vee}\phi)|_Y
\in\Gamma\bigl(Y,\Omega^1_{W/P^\vee}\otimes\mathcal O_Y(1,1)\bigr),
\]
其中\ednote{此处有难辨内容，请询问 A.G.。〔底本脚注37。〕}\textsuperscript{,}\transnote{右端的 $\Omega^1_{X/S}$ 均理解为拉回。截面取相对微分后限制在其零概形上，故更换线丛的局部平凡化不影响其定义；这说明上述 $d\phi$ 记号的确切含义。}
\[
\Omega^1_{W/P^\vee}\otimes\mathcal O_Y(1,1)
=\Omega^1_{X/S}\otimes\mathcal O_Y(1,1).
\]
我们将给出这个截面的另一种解释，结论由此立即得出。为此，考虑 $W$ 上，或更一般地，任意 $W$-预概形 $Z$ 上的如下正合列图。为免图形过宽，图中使用缩写
\[
\begin{aligned}
 A&=f^*\Omega^1_{P/S}\otimes\mathcal O_Z(1,0),& C&=N^\vee\otimes\mathcal O_Z(1,0),\\
 D&=\Omega^1_{X/S}\otimes\mathcal O_Z(1,0),& B&=E\otimes\mathcal O_Z,\\
 L&=\mathcal O_Z(0,-1),& M&=\mathcal O_Z(1,0),\\
 Q&=T_{P^\vee/S}\otimes\mathcal O_Z(0,-1).&&
\end{aligned}
\]
所有模都已拉回到 $Z$；$T_{P^\vee/S}=(\Omega^1_{P^\vee/S})^\vee$。横列及两条竖列正合，箭头方向均如下：\transnote{这是原图的校订重排，不是新增数学结论。原图左列顶端为 $\Omega^1_{X/Y}\otimes\mathcal O_Z(1,0)$，但由(a)拉回只能得到 $\Omega^1_{X/S}$。中列顶端原图记作难辨的 $G_{P^\vee/S}\otimes\mathcal O_Z(0,-1)$；由对偶Euler列确定为 $T_{P^\vee/S}\otimes\mathcal O_Z(0,-1)$，这里不声称读出了那个难辨字母。原图用类似比例号的符号表示复合，此处译校记为 $\alpha$。缩写不改变箭头的源与靶。}
\[
\begin{tikzcd}[row sep=small,column sep=large]
 &0&0&\\
 &D\arrow[u]&Q\arrow[u]&\\
 0\arrow[r]&A\arrow[u]\arrow[r]&B\arrow[u]\arrow[r]&M\arrow[r]&0\\
 &C\arrow[u]&L\arrow[u]\arrow[ur,bend right=15,"\alpha"']&\\
 &0\arrow[u]&0\arrow[u]&
\end{tikzcd}
\]
〔原编者向 A.G. 提问：上面的 $G$ 实际是一个难辨的字母，后接 $P^\vee/S$；这是什么？〕\ednote{请询问 A.G.。〔底本脚注38。〕}

第一条竖列由(a)拉回并张量 $\mathcal O_Z(1,0)$ 得出；横列由(b)拉回得出；第二条竖列由 $P^\vee$ 上类似的正合列 $(b^\vee)$ 的对偶列拉回得出，其中 $(b^\vee)$ 是把 $E$ 换成 $E^\vee$ 所得到的列。\transnote{底本将这些步骤略称“张量 $\mathcal O_Z$”，此处译校明示拉回。三条列均为局部分裂的局部自由模正合列，所以任意 $Z$ 上的拉回仍正合。由 $0\to\Omega^1_{P^\vee/S}(1)\to E^\vee_{P^\vee}\to\mathcal O_{P^\vee}(1)\to0$ 对偶，正得到 $0\to\mathcal O_{P^\vee}(-1)\to E_{P^\vee}\to T_{P^\vee/S}(-1)\to0$。}
由 $Y$ 的定义，$Z\to W$ 经 $Y$ 分解，当且仅当图中的复合态射 $\alpha:L\to M$ 为零，即当且仅当可以找到分解\transnote{底本此处源项印成 $\mathcal O_Z(0,01)$；由中间竖列及 $L\to B$ 的源确定，应为 $\mathcal O_Z(0,-1)$。}
\[
\beta:\mathcal O_Z(0,-1)\longrightarrow
f^*\Omega^1_{P/S}\otimes\mathcal O_Z(1,0).
\]
在这种情形下，可将它与
\[
f^*\Omega^1_{P/S}\otimes\mathcal O_Z(1,0)
\longrightarrow\Omega^1_{X/S}\otimes\mathcal O_Z(1,0)
\]
复合。把这一复合同态视作 $\Omega^1_{X/S}\otimes\mathcal O_Z(1,1)$ 的截面，得到的正是刚才引入的 $\Psi$ 的拉回。\ednote{Blass：请检查这里的字母是否正确。〔底本脚注39。〕}\textsuperscript{,}\transnote{底本写 $\Omega^1_{X/Y}\otimes\mathcal O_Z(1,1)$ 作为此同态的靶，混淆了同态与相应截面的扭转。正确同态为 $L\to D$；用 $\operatorname{Hom}(L,D)=\Gamma(Z,D\otimes L^{-1})$ 转成截面后才出现 $(1,1)$。$\psi$ 与 $\Psi$ 在此处译校统一作 $\Psi$。}
（验证应当基本上是机械性的。）

它为零，当且仅当 $Z$ 位于 $V(\Psi)$ 之上，这正是 $V(\Psi)$ 的定义！另一方面，这也意味着 $\beta$ 经由 $N^\vee\otimes\mathcal O_Z(1,0)$ 分解，即 $E\otimes\mathcal O_Z$ 的子模 $\mathcal O_Z(0,-1)$ 包含在子模 $N^\vee\otimes\mathcal O_Z(1,0)$ 中；这显然也就意味着 $Z\to W$ 经由
\[
\mathbf P(N(-1))\ \subset\ \mathbf P(E_X^\vee)
\]
这个子预概形分解。由此完成了所宣布的证明。\transnote{底本最后写作 $\mathbf P(\nu_{x/P}(1))$。由 $N^\vee(1)\subset E_X$ 对偶得到 $E_X^\vee\twoheadrightarrow N(-1)$，故相应闭浸入必须使用 $N(-1)$。作为抽象 $X$-概形，$\mathbf P(N(1))$ 和 $\mathbf P(N(-1))$ 都典范同构于 $\mathbf P(N)$；但用来定义本处闭浸入的商模及扭转必须写正确。}

在这个已经使我费了不少力气的博学句法练习之前，其实可以先指出：若 $S=\Spec(k)$，从集合论的角度看，非空的 $Y^{\mathrm{sing}}$ 的维数为 $r-1$，而 $P^\vee$ 的维数为 $r$，所以 $Y^{\mathrm{sing}}$ 在 $P^\vee$ 中的像的闭包余维至少为一，从而再次得到（2.12）。\transnote{底本说“像的余维至少一”，未明说闭包；这里按通常的像之维数解释写明闭包。空情形中结论自然成立，但不将空集维数写为 $r-1$。}值得注意的是，这个论证实质上并不不同于（2.12）中使用的论证。我们还注意到，这个集合的余维往往确实为一，参见下文。

因此，一般不能找到每个成员都光滑的超平面截面“线性束”。不过，我们将看到，往往仍能找到由没有任何超奇异点的超平面截面组成的线性束，因为在最常见的情形中，$Y^{\mathrm{sup\,sing}}$ 在 $P^\vee$ 中的像具有余维至少二。\transnote{底本此处写“余维二”。这里按下文所使用的余维下界表述，不声称必然恰等于二；例如异常集可以为空。这仍是原稿对后文结果的预告，不是此处译校已经证明了所有情形下的超奇异点判据。超奇异点术语沿用原稿后文定义。}
首先回顾当前情形中那些本质上属于微分性质的要点：


% DOT_IMPORT 04 8.7及相邻残段_校订续译工程

\egas{8.7}{定理的可读部分}
\noindent\textbf{(a)} 在目前的情形下，子预概形 ${Y^{\mathrm{sing}}}$（定义见第16号或第16节）正是 $\mathbf P(N)$，按前面说明的方式将它视为 $Y$ 的子概形。

\noindent\textbf{(b)} 预概形 ${Y^{\mathrm{sup\,sing}}}$ 的底空间（参见第16号或第16节），正是下述态射的分歧点集合：
\[
q:{Y^{\mathrm{sing}}}=\mathbf P(N)\longrightarrow P^\vee.
\]
这里源与靶均相对于 $S$ 光滑，非空的相关部分的相对维数分别为 $r-1$ 和 $r$。也就是说，这个态射在点 $y$ 处非分歧（参见定义），当且仅当 $y$ 对于 $\phi_\xi$ 是几何奇异而非超奇异零点（即本处二次型非退化的奇异零点），其中 $\xi$ 是 $y$ 在 $P^\vee$ 中的像。\transnote{原句用的是“geometrically an ordinary singular point”，即“几何普通奇异点”。底本第73--74页作者另行说明：此前把非退化二次奇异零点叫作“ordinary”与惯用术语冲突，应改称“singular non-supersingular zero”或“singular quadratic non-degenerate zero”。故此处译校采用“非超奇异零点”，限定为这里极化双线性型非退化的二次条件，不把它无条件等同于惯用的普通双点。特征二时奇数个变量的极化型总退化，而 $u^2+vw=0$ 的原点仍是惯用的普通二次奇异点；这一术语差异不能忽略。此处引用后文只为术语校勘，不计作已连续翻译后文。}

\bigskip
\noindent\textbf{(c)〔待核残段〕} 假设 $S=\Spec(k)$，且 $y\in{Y^{\mathrm{sing}}}=\mathbf P(N)$ 为一个 $k$-有理点。令〔字迹不清〕\ednote{请询问 A.G.。〔底本脚注40。〕}和 $\xi$ 分别为它在 $X(k)$ 和 $P^\vee(k)$ 中的投影；考虑 $P^\vee$ 中作为切映射之“像”的线性子簇 $H^1$……

\noindent〔原句在此与下一段内容混接，没有给出可辨认的完整切空间结论。〕\transnote{原文从“linear subvariety $H^1$ … image of the tangent map”直接连到“of the closure of its … image in $P^\vee$”，其语法、定义和论证均有断裂；并附法文“son image”（它的像）。这里的省略号明确表示底本连接残缺，不代表主动略去一个已读懂的证明。此处译校没有补造常见的对偶切空间定理。}

\subsection*{紧随（8.7）(c)的条件列表与论证残段}
\noindent\textbf{〔编号及假设待核〕} 取其在 $P^\vee$ 中的像的闭包，赋予诱导的约化结构，并考虑诱导态射
\[
g:{Y^{\mathrm{sing}}}\longrightarrow T,
\]
它是整预概形之间的支配态射。\transnote{底本此处未完整引入 $T$，也未明确“其”指代的映射。依该句自身及后续记号，$T$ 是所讨论像闭包的名称；但此处译校不据此认定遗失假设已经恢复。前文 $X/S$ 光滑并不自动保证 $X$ 或 ${Y^{\mathrm{sing}}}$ 整，故必须把“整预概形”视为本段另行要求的条件。后文回指（8.8），本页却没有其标题，此处译校不擅自补号。}

原稿随后列出下列条件。为清楚起见，$L/K$ 表示由这个支配态射诱导的函数域扩张，$K=k(T)$、$L=k({Y^{\mathrm{sing}}})$。\transnote{函数域在整性及非空条件下有定义。底本仅说“the field extension $L/K$ defined by $g$”，此处译校显写其含义，不为其他缺失假设作补充。}
\begin{enumerate}
\item[(i)] 态射 $g$ 泛平展，即至少在一点平展；等价地，在 ${Y^{\mathrm{sing}}}$ 的泛点平展，也就是在该泛点非分歧。
\item[(i bis)] $L/K$ 为有限可分扩张。
\item[(i ter)] $g$ 是双有理态射，即 $L/K$ 为平凡扩张。
\item[(ii)] 从集合论意义说，${Y^{\mathrm{sing}}}\ne{Y^{\mathrm{sup\,sing}}}$。
\item[(ii bis)] 存在 $x\in X(k)$，以及在 $x$ 处与 $X_k$ 相切但在该点不密切的超平面 $H$；这里“不密切”的确切含义是：对定义 $H$ 的 $\mathcal O_{X_k}(1)$ 的截面而言，$x$ 不是超奇异点。
\end{enumerate}

在上述整性、支配性及有限型背景下，可以确认的关系是
\[
\text{(i)}\ \Longleftrightarrow\ \text{(i bis)},
\qquad
\text{(i ter)}\ \Longrightarrow\ \text{(i)}.
\]
原稿要求的全组等价关系，在现存假设下不能成立，须等待缺文核查。\transnote{原句译为：“条件(i)与(ii)(bis)是等价的”，但随后实际列出五个条件，证明还要求证明(i)$\Rightarrow$(i ter)。后一蕴涵在仅有“$f$非分歧”时为假。取 $k=\mathbb C$、$X=\mathbb G_m$、$f(t)=[1:t^2:t^4]$。此映射是二次平展覆盖后接到圆锥曲线开子集的同构，再接该子集到 $\mathbf P^2$ 的浸入，故非分歧；${Y^{\mathrm{sing}}}\simeq X$，而 $g(t)=[t^4:-2t^2:1]$。其像闭包为光滑对偶圆锥曲线，$L/K=\mathbb C(t)/\mathbb C(t^2)$ 次数为二且可分，因此(i)成立而(i ter)不成立。必须补充能排除此类多重覆盖的适当假设；例如 $f$ 为浸入可排除此反例，但此处译校不恢复缺失的原假设。}

此外，在未假设 $k$ 代数闭的现存文字下，(ii)也不能直接替换为存在 $k$-有理点的(ii bis)。\transnote{非空的几何开集未必有 $k$-有理点。例如取 $k=\mathbb R$，$X\subset\mathbf P^2$ 为光滑圆锥曲线 $x_0^2+x_1^2+x_2^2=0$。它没有实点，故(ii bis)失败；但在 $\mathbb C$ 上每条切线的接触均是非退化二重接触，故 ${Y^{\mathrm{sup\,sing}}}=\varnothing$ 而 ${Y^{\mathrm{sing}}}\ne\varnothing$，(ii)成立。可通过使用几何点或补入代数闭条件修正这一关系，但原稿缺文未能确定作者究竟采用哪种表述。}

当 ${Y^{\mathrm{sing}}}$ 非空、整、维数 $r-1$，且(ii)成立时，作为分歧闭集的 ${Y^{\mathrm{sup\,sing}}}$ 是真闭子集，因而若其非空，便有
\[
\dim{Y^{\mathrm{sup\,sing}}}\leq r-2.
\]
所以它在 $P^\vee$ 中的像的闭包余维至少为二。\transnote{底本此句印成：“这些条件蕴含 ${Y^{\mathrm{sup\,sing}}}\ne\varnothing$〔法文难辨，请询问 A.G.〕$\dim{Y^{\mathrm{sup\,sing}}}\leq r-2$，从而其像余维至少二。”非空性不能由前述条件推出；例如在代数闭特征零域上的光滑非退化圆锥曲线没有超奇异切截面。因此正文仅保留经维数论可核实的条件式：${Y^{\mathrm{sup\,sing}}}$ 为空，或其维数至多为 $r-2$。}
若(i bis)成立，还得到
\[
\text{(iii)}\qquad\dim T=r-1,
\]
即 $T$ 在 $P^\vee$ 中余维为一。

\bigskip
\noindent\textbf{证明的现存前段。}
(i)与(i bis)的等价是直接的。至于它与(ii)的等价，需要把 $g$ 到像闭包的非分歧性与 $q:{Y^{\mathrm{sing}}}\to P^\vee$ 的非分歧性在泛点比较；在本段的整性和有限型背景下，可由（8.7）(b)作此比较。\transnote{底本原句：“它与(ii)的等价是（8.7）b)的一个直接推论；最后(ii)与(ii bis)的等价，实际上就是 $H^{\mathrm{sup\,sing}}$ 的定义。”后一段中的 $H^{\mathrm{sup\,sing}}$ 疑为 $Y^{\mathrm{sup\,sing}}$，且其有理点断言需要上面说明的基域条件，故不将其作为无条件的正确证明保留在正文。对前一比较，$T\hookrightarrow P^\vee$ 为闭浸入；若 $g$ 泛平展则复合泛非分歧；反之复合泛非分歧蕴含 $g$ 泛非分歧，再由函数域有限可分性得到泛平展。}
显然，(i ter)蕴涵(i)。其逆命题在现存假设下为假，所以不能按原稿所写继续完成一个无条件的证明。\transnote{以下完整保存原稿的后续证明文字，而不把它作为有效证明：“还须证明(i)$\Rightarrow$(i ter)。显然可以假设 $K$ 是代数闭的；考虑到假设(i)，归结为证明存在一个非空开集 $U$，使得对每个 $\xi\in U(K)$，在 $\xi$ 上方恰有一个 ${Y^{\mathrm{sing}}}(K)$ 的点。这将由（8.7）c)推出；后者更精确地蕴含下面的断言。”原文将先前函数域 $K=k(T)$ 与基域混用，不能把一个一般函数域无说明地“假设为代数闭”；即便将此理解为对基域作代数闭扩张，二次覆盖反例仍不消失。它所援引的(c)切空间结论在此底本中残缺。因此不能据这段话填补唯一性证明。底本末处(i er)显为(i ter)的排印错，本脚注按所列条件编号书写。}


% DOT_IMPORT 05 8.9至8.11_校订续译工程

\egas{8.9}{推论：现存条件下的可核部分}
\noindent\textbf{译校记号约定。}$Q={Y^{\mathrm{sing}}}$，$R={Y^{\mathrm{sup\,sing}}}$，$q:Q\to P^\vee$；$T$ 为像闭包的约化结构，$g:Q\to T$，$U=T_{\mathrm{sm}/k}$。在使用函数域时假设 $Q,T$ 非空且整，并沿用有限型背景；$K=k(T)$、$L=k(Q)$，$r=\dim P^\vee$，$\dim Q=r-1$。$R$ 按（8.7）(b)为 $q$ 的分歧闭集。$Q_U=g^{-1}(U)$。这些条件在下文使用时明示，不声称已恢复原稿缺失的整套假设。
假设（8.8）的条件(i)成立，即 $g$ 泛平展，令 $U$ 为 $T$ 中相对于 $k$ 光滑的点所成的开集。可以确认 $U\ne\varnothing$。至于 $Q_U\to U$ 是开浸入的结论，在现存假设下不成立；原稿完整断言及固有情形的后续结论见脚注，均不能仅据泛平展性采纳。\transnote{原推论译文：“假设（8.8）的条件(i)成立，令 $U$ 为 $T$ 中相对于 $k$ 光滑的点所成的开集。则 $U\ne0$，${Y^{\mathrm{sing}}}/U\to U$ 为开浸入，因而 ${Y^{\mathrm{sing}}}/U$ 不含 ${Y^{\mathrm{sup\,sing}}}$ 的点。若 $X$ 在 $K$ 上固有，则 $g:{Y^{\mathrm{sing}}}\to T$ 满射，从而 ${Y^{\mathrm{sing}}}/U\to U$ 在 $K$ 上固有，所以 $g:{Y^{\mathrm{sing}}}\to T$ 满射，因此 ${Y^{\mathrm{sing}}}/U\to U$ 是同构，且 $U$ 是具有后一性质的最大开集。”这里两个 $K$ 与前文函数域记号冲突；若讨论 $X$ 的基域，应为 $k$。此外，$U$ 上的态射固有应指相对于 $U$，不能推出其源在 $k$ 上固有。更重要的是，取 $k=\mathbb C$、$X=\mathbb G_m$、$f(t)=[1:t^2:t^4]$，则 $f$ 非分歧，$Q\simeq X$，$g(t)=[t^4:-2t^2:1]$ 到光滑对偶圆锥曲线泛平展、次数二；此时 $U=T$，$Q_U\to U$ 不是开浸入。故此处译校不把这条主结论作为正确正文。反例只针对现存假设，不能否定未知完整假设下的原定理。}

如果另已证明 $Q_U\to U$ 为开浸入，并且该态射固有且满射，那么它必为同构。这是原稿后半段可以单独核实的条件式，并不填补其前半段的缺口。\transnote{固有的开浸入，其像同时开且闭；再加满射便是同构。此论证没有证明当前 $g$ 满足开浸入假设。若 $Q$ 在 $k$ 上固有而 $T$ 在 $k$ 上分离，则 $g$ 固有；其支配性与固有性使像为全 $T$。这是正确的固有性传递方向。}

\subsection*{证明与缺口记录}
首先，由 $g$ 支配且泛平展，可以找到 $T$ 的一个非空开子集 $V$，使 $g^{-1}(V)\to V$ 平展且满射；于是 $V$ 在 $k$ 上光滑，从而 $U$ 非空。\transnote{原文把“泛平展”重复写了两次，并混用 $k$ 与 $K$；此处只去掉重复，统一基域为 $k$。这里的 $Q$ 在 $k$ 上光滑，光滑性沿平展满覆盖下降；泛平展且泛有限使适当缩小 $T$ 后可除去非平展部分的像，再利用平展像开性得到满射。}

对于 $\xi\in V(k)$ 及其上方的 $y\in Q(k)$，原稿试图利用（8.7）(c)从 $T$ 在 $\xi$ 处的切空间唯一恢复 $x$，进而恢复 $y$。这一唯一性推断没有得到现存假设的支持。\transnote{原证明对应段落完整译为：“若 $\xi\in V(K)$，且 $y$ 是其上方 ${Y^{\mathrm{sing}}}(k)$ 的一点，则沿用（8.7）c)的记号，空间 $H^0$ 正是 $T$ 在 $\xi$ 处的切空间；如我们在那里所指出的，$y$ 投影到 $X(k)$ 的点 $x$ 由与 $H^0$ 正交的点确定，因而唯一确定；又因为 ${Y^{\mathrm{sing}}}\subset X\times P^\vee$，所以 $y$ 也唯一确定。这已证明 $g$ 是双有理的，因为它泛平展且泛根性。”但(c)在底本中残缺，$H^0$ 也未完整定义。即使切空间确定了环境射影空间中的 $f(x)$，一般非分歧态射仍可能在此点上方有多个 $x$；二次覆盖反例正是如此。故不能从这段推断双有理性。$K$ 的大小写和 $H^0/H^1$ 差异均是底本原状。}

接下来构造逆映射的段落同样依赖该唯一性；因此此处译校保存其完整可辨文字作为待核原证明，不把它改造为已经建立的逆态射。\transnote{原文续译：“另一方面，考虑态射 $\psi$（其定义形式显然）：它给每个 $\xi\in U(K)^0$ 对应 $P$ 中与 $U$ 在 $\xi$ 处切空间正交的唯一点 $x=\psi(\xi)$；在 $V$ 上，它与复合 $V\to{Y^{\mathrm{sing}}}|_V\to X$ 一致，第二个箭头为投影。因此，令 $h=(\psi,\mathrm{id}):U\to P\times T$〔字迹不清〕，$g_1=g|_{g^{-1}(U)}:g^{-1}(U)\to U$，则复合 $hg_1:g^{-1}(U)\to P\times T$ 恰为典范包含，因为在 $g^{-1}(V)\xrightarrow{\sim}V$ 上它是如此。于是 $h$ 经由 $P\times T$ 中 $Y^1$ 的概形论闭包 $Y^1$ 分解；所以 $Y^1$（它在上述闭包中开）在 $h$ 下的逆像是 $U$ 的开子集，记为 $U^1$。由于 $hg_1=\text{包含}$，立即可见 $g^1$ 诱导同构 $g^{-1}(u)\xrightarrow{\sim}W$ 是同构；由 $Y^1$ 光滑可知 $W$ 光滑，所以 $W\subset U$。这就证明了（8.9）。”底本在闭包与原子概形处均印作 $Y^1$，并有 $U^1,W,u,g^1$ 混用；这些对象尚不能唯一辨认，未据猜测修复。原脚注41为“请询问 A.G.”，已在此保留。即便处理记号混乱，此段仍需要前面未证明的唯一性与浸入类条件。}

\subsection*{关于（8.8）末尾维数断言的补充}
若 $Q$ 非空、整且维数为 $r-1$，$R$ 为真闭子集，则 $R$ 为空，或者 $\dim R\leq r-2$。要把不等式加强为非空时恰等于 $r-2$，还须证明这里的 $R$ 确由一个非零可逆模截面定义；此处译校不据缺字补出这一识别。\transnote{原文：“（8.8）最后的断言 ${Y^{\mathrm{sup\,sing}}}=\varnothing$ 或 $\dim{Y^{\mathrm{sup\,sing}}}=r-2$，以及 $\dim T=r-1$，都是平凡的：第一个来自 ${Y^{\mathrm{sing}}}$ 不可约且维数为 $r$，以及 ${Y^{\mathrm{sing}}}$ 或 ${Y^{\mathrm{sup\,sing}}}$〔字迹不清〕是可逆模某截面 $D$ 的消失所定义的这一事实。”其中 $\dim{Y^{\mathrm{sing}}}=r$ 与前面射影法丛的计算及本句结尾均冲突，应为 $r-1$。底本原脚注42为“请询问 A.G.”。此处译校不将有待辨认的零点对象和截面认定为已经恢复。}
第二个维数等式来自 $L/K$ 有限，因而
\[
\operatorname{trdeg}_k L=\operatorname{trdeg}_k K,
\qquad\dim T=\dim Q=r-1.
\]

\egas{8.10}{备注与集合差的订正}
沿用（8.9）的记号，在已另行获得该推论所需前提和开浸入结论时，有
\[
g^{-1}(U)\subset Q-R.
\]
原稿指出，即使 $X$ 闭于 $P$，这一包含也未必是等号。令 $E_{\mathrm{et}}$ 为 $g$ 的平展点集，则可独立确定：“非分歧而不平展”的点集是
\[
(Q-R)\setminus E_{\mathrm{et}}.
\]
只有另已证明 $E_{\mathrm{et}}=g^{-1}(U)$ 时，才能把上式写成 $(Q-R)\setminus g^{-1}(U)$。\transnote{底本说“${Y^{\mathrm{sing}}}-{Y^{\mathrm{sup\,sing}}}$ 是 $g$ 非分歧但不平展的点集”，漏掉了集合差。按（8.7）(b)及 $T\hookrightarrow P^\vee$ 为闭浸入，$Q-R$ 是 $g$ 的非分歧点集，其中还包含平展点，必须去掉平展点集。原稿把 $g^{-1}(U)$ 认作平展域，仍依赖待核的（8.9）背景，此处译校不将两者直接等同。}

在本段有限型、整、支配、同维的背景下，非分歧而非平展点的像不是几何单支点，因而也不是几何正规点。\transnote{原文为“这还蕴含 $g(y)^0$ 是 $T$ 中一个不几何正规、甚至不几何单支的点”；上标 $0$ 含义不清，正文只写像点 $g(y)$。独立核验的理由是：非分歧态射在靶的平展局部可写成闭浸入到一个平展邻域；若靶在该点几何单支，则相应平展局部的约化结构只有一个局部分支。在这里靶整、局部平展代数约化，因而相应局部环整。缩小到该点的整平展邻域后，其维数与源均等于 $\dim T$；闭浸入的像因同维而占满这个整邻域，故其定义理想为零，故此局部闭浸入为同构，原态射在该点平展，矛盾。这个结论使用了本段的支配、同维条件，不能仅凭“非分歧但不平展”推出，也不证明（8.9）的开浸入断言。}

从几何上说，这对应于如下现象：一个超平面在 $x\in X(k)$ 处与 $X$ 相切但不密切，同时还可能存在另一点 $x^1\in X(k)$，使同一超平面在 $x^1$ 处相切。\transnote{底本句末写“at $x$ [or $x^1$ illegible]”，译为“在 $x$ 处〔或者是 $x^1$，字迹不清〕”；由该句明确讨论另一接触点的语义，正文用 $x^1$。原脚注43为“字迹不清”，保留于此。不密切仍按前文解释为相应二次型非退化，而非随意替换为其他相切阶数定义。}
原稿举例说，可以令 $P$ 维数为二、$X$ 为次数至少四的非奇异曲线，并称在任意特征下都有显然的例子。〔这里是在转述原稿的例子预告，此处译校未把它加强为“所有这样的曲线都有双切线”。〕〔附注：$X$ 的“双切线”对应于“对偶曲线”的双点。〕

\egas{8.11}{特征零的推论}
假设 $k$ 的特征为零，保持上文 $Q$ 非空且整、光滑、维数 $r-1$ 等背景。

\noindent\textbf{(a)} $R={Y^{\mathrm{sup\,sing}}}$ 在 $P^\vee$ 中的像的闭包余维至少为二。

\noindent\textbf{(b)} 条件(iii)，即 $\dim T=r-1$，等价于 $L/K$ 有限，亦等价于(i bis)、(i)及(ii)。这里(ii)为 $Q\ne R$。因而其否定 $Q=R$ 等价于 $Q$ 在 $P^\vee$ 中的像的闭包余维至少为二。\transnote{原文说“（8.8）的条件(iii)与其他条件等价”，继而说明其他条件的否定，即 ${Y^{\mathrm{sing}}}={Y^{\mathrm{sup\,sing}}}$，意味着 ${Y^{\mathrm{sing}}}$〔或 ${Y^{\mathrm{sup\,sing}}}$，字迹不清〕在 $P^\vee$ 中的像余维至少二。若“其他条件”包含(i ter)的双有理性，则这一说法即便在特征零也被二次平展覆盖反例否定；若包含(ii bis)的 $k$-有理点存在，还缺少代数闭等条件。此处译校只保留可由现存整性、维数、有限型和特征零条件证明的等价关系。最后的难辨对象在 $Q=R$ 时无底集差别，故正文写 $Q$，并在此留存原文两种读法。}

\begin{proof}
考虑到前述维数结论，(b)蕴含(a)。由维数理论，(iii)意味着 $L/K$ 为有限扩张；也可以把它写成（8.8）中的条件“(iii bis)”。而特征零时 $L/K$ 总是可分的，所以得到（8.8）的(i bis)。这就是原稿给出的论证。

为避免依赖缺失的全组等价断言，补充说明如下〔此处译校补证〕：$T$ 是 $q$ 的像闭包，故 $\dim T\leq r-1$。若 $\dim T\leq r-2$，则 $R$ 的像闭包包含于 $T$，(a)立即成立；且 $q$ 不可能在任何点非分歧，否则整性使其在泛点非分歧，从而 $L/K$ 有限，迫使两者维数相等。故此时 $R=Q$。

若 $\dim T=r-1$，则 $L/K$ 有限，特征零保证可分，因此 $g$ 泛平展，$q$ 泛非分歧。由（8.7）(b)，$R$ 为 $Q$ 的真闭子集，于是 $R$ 为空或维数至多为 $r-2$，其像的闭包余维至少为二。反之，$Q\ne R$ 使非分歧开集非空并含泛点，故 $L/K$ 有限可分，得到(iii)。这证明了正文中明确列出的等价关系，也直接证明(a)，没有使用未经证实的双有理性。
\end{proof}

% DOT_IMPORT 06 8.12_校订续译工程
\egas{8.12}{备注}
作者从几何上将断言(a)解释如下：对于一个充分一般的超平面截面线性束，束中的每个成员或者光滑，或者其几何奇点都是普通二重点。
\transnote{此处(a)、(b)指紧接前面的8.11。原文还声称：“事实上，立即可见，因而还可以把(a)说得稍精确些：至多有一个这样的几何奇点。”在目前缺失假设的底本中，不能直接采纳这一加强。前文译校所给 $k=\mathbb C$、$X=\mathbb G_m$、$f(t)=[1:t^2:t^4]$ 的非分歧二次覆盖，对 $a\ne0$，相切超平面 $a^2Z_0-2aZ_1+Z_2=0$ 的拉回为 $(t^2-a)^2$，在 $t=\pm\sqrt a$ 两点各产生一个二重零点，故不能仅从现存非分歧假设推出“至多一个”。这里保留作者原断言以供对照，但不声称已验证它。普通二重点与8.8所讨论非退化性的具体对应也需按完整假设核查。}

断言(b)的基本含义是：如果 ${Y^{\mathrm{sing}}}={Y^{\mathrm{sup\,sing}}}$，那么，对于一个充分一般的超平面截面线性束，束中所有成员都是光滑的。\transnote{此句是对8.11(b)的解释，沿用该推论的特征零背景及相关前提，并非对任意特征的无条件断言。原文下一句“无论是否特征零”修饰的是对这种情形的“例外”评价，不应倒读为已证明任意特征的结论。}
上述等式可在解析意义下表述为：$Y^1$ 上某个可逆模的一个截面 $D$ 消失；底本将该可逆模印作\transnote{底本这里同时使用 $Y^1$ 和 $Y'$，且未明确标出 $\omega_{X/k}$ 到相应空间的拉回。此公式按可辨印文保留，记号与定义域待核，不把它视为已恢复的精确线丛公式。}
\[
\omega_{X/k}^{\otimes2}\otimes\mathcal O_{Y'}(1,1).
\]


这第二种情形，无论我们是否处于特征零，都应当完全视为例外情形。〔页顶手迹残段：“簇 $L$ 在 $1\ldots T=L$”；字迹不清。〕\footnote{原脚注44：“字迹不清”。}
用古典语言来说，如果没有弄错，就是说：相对于所考虑的射影浸入，$X$ 是“直纹的”（ruled）。\transnote{此为作者此处暂用的术语，不能直接等同于通常的直纹簇定义。底本后文8.17又明确表示这里的用法与传统用法不一致；此处译校不替作者作未经确认的重新定义。}
〔而且，如果我们愿意，〕由8.5及其推论，这里已有把这一术语明确化并加以论证所需的一切；如果你有兴致，也可把它同〔法文：\emph{la taupe}〕联系起来。\transnote{底本保留法文词组而未解释其指代；本译不猜作某个数学定理或术语。}
例如，若 $\dim X=1$，这意味着 $X$ 是一条直线〔字迹不清〕\footnote{原脚注45：“字迹不清”。}，$x\in X(k)$，所以 $T$ 包含〔缺文〕\footnote{原脚注46：“字迹不清”。}。
\transnote{“直线”句及随后包含关系在底本中断裂，无法据此确定所需浸入、闭性或其他限制；此处只保存可辨残句，不能将其作为完整已证命题使用。}

(b) 若特征为 $p>0$，通常应当给出这样的例子：$\dim P=2$，$X$ 为非奇异代数曲线，8.6的条件不满足，即 ${Y^{\mathrm{sing}}}={Y^{\mathrm{sup\,sing}}}$，但仍有 $\dim T=r-1$；也就是说，给出 $L/K$ 为有限不可分扩张的例子。我太懒了，不愿构造这些例子，但我不怀疑它们存在。\footnote{原脚注47：“把它做出来。——Blass”。}\transnote{原文此处确实写“8.6”，并重新出现“(b)”编号；两者均保留，引用关系待核。作者表达的是对例子存在的确信，不是已完成的构造或证明。}

\subsection*{原稿编排备忘}
在(a)中加一条脚注，指向后面的编号小节或段落；在那里我们证明：如果出现这一例外的“直纹”情形，那么对射影浸入作一个简单的修改，就能重新回到8.8的“一般”情形。\transnote{这是作者对后文的编排指示，并非此处译校已经完成该修改或核验该证明的声明。}

本节从8.6到这里的部分，无疑可以独立成一个具有微分性质的小节；而本编号小节的开头，应当与接下来的部分合并成一个讨论例外〔超平面？〕\footnote{原脚注48：“请询问 A.G.”。}之维数的小节。我这里只用了 ${Y^{\mathrm{sing}}}$ 的维数是 $r-1$ 这一事实。


% DOT_IMPORT 07 8.13_校订续译工程
\egas{8.13}{命题（原稿陈述及待核处）}
仍假设 $f\colon X\to P$ 非分歧，且 $X$ 没有孤立点。假设 $X$ 在几何意义下满足 $(R_k)$。

令 $Z_k$ 为如下点集在 $P^\vee$ 中的补集：这些点 $\xi\in P^\vee$ 满足 $\phi_\xi$ 是 $X_{k(\xi)}$-正则的，且 $Y_\xi$ 满足几何条件 $E_k$。\transnote{底本此处确实印作 $E_k$，而前句为 $(R_k)$。没有可靠依据将二者擅自统一，故保留 $E_k$ 并标待核。这里下标 $k$ 与基域符号 $k$ 沿用底本；$k(\xi)$ 表示点 $\xi$ 的剩余域。}则：
\begin{enumerate}
\item[(a)] 为使 $Z_{k-1}$ 在 $P^\vee$ 中的余维至少为 $2$，只须 $X'$ 的每个维数不超过 $k$ 的不可约分支 $x'_i$，相对于 $f$ 都是“直纹的”。\transnote{底本分支记号印作小写 $x'_i$，此处不猜改为 $X'_i$。“直纹的”（ruled）沿用作者的临时用语；8.17明确提醒此用法与传统术语不一致，不能径按通常直纹簇定义理解。}
\item[(b)] 为使 $Z_k$ 在 $P^\vee$ 中的余维至少为 $2$，只须每个维数不超过 $k-1$ 的不可约分支 $X_i$，都由 $X$ 的光滑点组成，并且是“直纹的”。\transnote{底本在“维数不超过 $k-1$”后插有“询问 A.G.；字迹不清”。该条件带有原稿疑问，尚未由原始手稿恢复或独立核实；本译保存可辨条件，不把它改写成已经审定的充分条件。}
\end{enumerate}


\subsection*{原稿证明段落}
事实上，只考虑使 $\phi_\xi$ 为 $X_{k(\xi)}$-正则的 $\xi$。对这样的每个 $\xi$，$Y_\xi$ 的几何奇点集是 $\operatorname{sing}(Y'_\xi)$ 与 $T$ 在 $Y_\xi$ 中的逆像 $T_\xi$ 的并。\transnote{底本首句的英语语序断裂：“for every $\xi$ geometrically singular the set of $Y_\xi$”；这里按随后的并集陈述整理语序，并不增加数学假设。$T$ 的指代必须沿本节上下文核定，不能误认成前面出现的任意同名集合。}
（注意：我们只考虑使 $\phi_\xi$ 为 $X_{k(\xi)}$-正则的 $\xi$；由8.2，这一限制无妨。）\transnote{底本此处误排为 $\phi_{x}i$，按命题开头同一正则性条件订正为 $\phi_\xi$。}
因此，这一奇点集在 $Y_\xi$ 中的余维等于下述两个余维的下确界：
\transnote{底本第一个余维式印作 $\operatorname{codim}(\operatorname{sing}(Y'),Y'_\xi)$，遗漏前文及后文均有的纤维下标 $\xi$；正文按其所讨论的纤维奇点集补足。此为记号一致性的排印订正，不是对余维等式适用前提的新增证明。}
\[
\operatorname{codim}\bigl(\operatorname{sing}(Y'_\xi),Y'_\xi\bigr),
\qquad
\operatorname{codim}(T_\xi,Y_\xi).
\]


现在只考虑使 $\operatorname{sing}(Y'_\xi)$ 有限的 $\xi$。（这一限制无妨；这使我们处于一个余维至少为 $2$ 的集合的补集中。）于是，$Y'_\xi$ 的几何奇点都是孤立的。原稿据此及8.4断言，结论容易得到。\transnote{原文末句为“The conclusion follows easily from this and from 8.4.”。底本未展开由这些事实推出(a)、(b)的逐步推导，此处译校不补造它；特别是 $E_k$ 记号、(b)条件及“直纹”含义仍待核，故不能将本段标作完整独立审定的证明。}

% DOT_IMPORT 08 8.14至8.15_校订续译工程
